\documentclass[UTF8,a4paper,12pt]{ctexart}
\usepackage{geometry}
\usepackage{booktabs}
\usepackage{hyperref}
\usepackage{xcolor}
\usepackage{listings}
\geometry{left=2.5cm,right=2.5cm,top=2.4cm,bottom=2.4cm}
\hypersetup{colorlinks=true,linkcolor=blue,urlcolor=blue}
\setlength{\parindent}{2em}
\setlength{\parskip}{0.3em}
\lstset{basicstyle=\ttfamily\small,breaklines=true,columns=fullflexible,
  frame=single,numbers=left,numberstyle=\tiny,showstringspaces=false,
  keywordstyle=\color{blue!55!black},commentstyle=\color{green!35!black}}

\title{预备工作实验报告\\了解编译器、LLVM IR 编程及 RISC-V 汇编编程}
\author{姓名：\underline{\hspace{3.5cm}}\quad 学号：\underline{\hspace{3.5cm}}\\
组员：\underline{\hspace{3.5cm}}\quad 班级：\underline{\hspace{3.5cm}}}
\date{\today}

\begin{document}
\maketitle

\begin{abstract}
本实验根据课程 PPT 中“了解编译器、LLVM IR 编程及汇编编程”的要求，以阶乘程序为统一实验对象。首先使用 GCC 和 Clang 分别生成预处理文件、LLVM IR、宿主机汇编、目标文件和最终可执行文件，并逐阶段分析源程序中的变量、循环、运算和函数调用如何被降低。随后手写与源程序等价的 LLVM IR，使用 SSA 形式和 \texttt{phi} 指令表示循环数据流；再选择 RISC-V 作为目标架构，手写 RV64 汇编主程序和只依赖 Linux 系统调用的最小 SysY 兼容运行时，经过汇编、链接后在 QEMU 中运行。进阶部分比较了 \texttt{-g}、\texttt{-O0}、\texttt{-O2} 和目标架构选项造成的中间代码差异，并使用边界输入和普通输入验证三种实现。实验结果表明，C、手写 LLVM IR 和手写 RISC-V 汇编在输入 0、1、5、7 时输出完全一致。
后续验证中，直接连接所提供的 x86 SysY 静态库完成 30 次检查，并从所提供的未修改源码重建 RISC-V Linux 运行库，完成 45 次检查；原自写运行时路线扩展为 39 次检查。原包内 RISC-V 静态库则因当前工具链及标准库环境不匹配而未能直接链接，报告如实区分该失败与源码重建成功。
另设计包含常量、数组、条件分支、循环和用户函数的 SysY 示例，手写对应 LLVM IR 与 RV64 汇编，在两个目标平台上完成 40 次结果检查。此报告针对课程的 5 分预备工作；后续实现词法分析器、语法分析器和完整编译器的模块不计入本次完成情况。
\end{abstract}

\noindent\textbf{关键词：} 编译过程；LLVM IR；SSA；RISC-V；SysY 运行时；交叉编译

\section{引言与实验任务}

课程 PPT 对本次预备工作的要求可以拆成以下五项。

\begin{enumerate}
  \item 以一个简单 C/C++ 程序为例，使用命令行选项得到预处理、编译、汇编和链接各阶段的输出，并分析中间产物与源程序的关系。
  \item 设计一个包含编译器所支持语言特征的 SysY 示例程序。
  \item 写出与示例程序语义等价的 LLVM IR 和 ARM 或 RISC-V 汇编程序，连接运行时并验证结果。
  \item 通过调试选项、优化选项、目标架构选项或修改测试输入，进一步观察输出变化。
  \item 按科技论文结构完成 PDF 报告，并说明真实的小组分工。PPT 中的 LaTeX 资料与模板是建议使用的参考，不将其误写为强制格式。
\end{enumerate}

本实验先选择阶乘程序观察完整编译链，再增加特性示例来覆盖阶乘没有涉及的条件分支、用户函数和数组。它们共同服务于《预备工作——了解你的编译器》PPT 所列的 5 分实验。课程总体要求文档还包含词法、语法、优化和目标代码生成等后续独立模块，这些不是此份预备实验报告声称完成的内容。

\section{实验环境、文件和分工}

实验在 Linux x86-64 主机上完成。实际工具为 GCC 13.3.0、Clang/LLVM 18.1.3、RISC-V GNU Binutils 2.42 和 QEMU 8.2.2。查看版本的命令为：

\begin{lstlisting}[language=bash]
gcc --version
clang --version
riscv64-unknown-elf-as --version
qemu-riscv64 --version
\end{lstlisting}

所有命令均在项目根目录执行。实验文件如下：

\begin{lstlisting}
src/factorial.sy                 SysY 风格源程序
src/factorial.c                  供 GCC/Clang 使用的 C 版本
src/features.sy                 多特征 SysY 示例
src/features.c                  同算法的 C 版本
src/sysy_runtime.h               getint/putint 函数声明
src/sysy_runtime.c               宿主机运行时实现
ir/factorial_manual.ll           手写 LLVM IR
ir/features_manual.ll            多特征手写 LLVM IR
asm/factorial_riscv64.s          手写 RISC-V 主程序
asm/features_riscv64.s           多特征手写 RISC-V 程序
asm/sysy_runtime_riscv64.s       手写 RISC-V 运行时
lib/                             提供的运行库、来源说明与校验值
tests/check_factorial.py          结果、输出字节和退出状态检查
tests/check_features.py           多特征分支测试
build/                           各阶段生成物和可执行文件
\end{lstlisting}

小组分工需要按实际情况填写，不能直接保留模板文字。依照 PPT，两位组员都应独立完成编译各阶段的观察，并分别记录自己的操作和分析；在等价实现部分，一人负责 LLVM IR，另一人负责 RISC-V 汇编。两人可以共享报告框架、示例和测试结果，但应各自提交报告。此处不预先认定任何组员已经完成其个人实验。

\section{实验总体路线}

第一条路线用于理解现有编译器：

\begin{lstlisting}
C 源程序 -> 预处理文件 -> LLVM IR/汇编 -> 目标文件 -> 链接 -> x86-64 程序
\end{lstlisting}

第二、三条路线用于理解中间表示和目标代码：

\begin{lstlisting}
SysY 算法 -> 手写 LLVM IR -> 与 C 运行时链接 -> x86-64 程序
SysY 算法 -> 手写 RISC-V 汇编 -> 与 RISC-V 运行时链接 -> RISC-V 程序
\end{lstlisting}

最后对三种实现运行同一组输入，并与独立的阶乘期望值比较。有限测试提供这些输入下的正确性证据，不能作为所有输入上的语义等价证明。步骤十三另外区分自写运行时与提供的 SysY 库，避免把二者混为同一项验证。

\section{步骤一：设计 SysY/C 示例程序}

\subsection{目的和做法}

本步对应“选择简单程序”和“设计覆盖语言特征的 SysY 程序”。程序读取整数 $n$，先令 \texttt{i=2}、\texttt{f=1}；当 \texttt{i<=n} 时执行 \texttt{f=f*i} 和 \texttt{i=i+1}；循环结束后输出 \texttt{f}。

文件 \path{src/factorial.sy} 的内容如下：

\begin{lstlisting}[language=C]
int main() {
    int i;
    int n;
    int f;
    n = getint();
    i = 2;
    f = 1;
    while (i <= n) {
        f = f * i;
        i = i + 1;
    }
    putint(f);
    return 0;
}
\end{lstlisting}

其中 \texttt{int} 覆盖整型变量和返回值，\texttt{getint/putint} 覆盖函数调用，\texttt{while} 覆盖循环控制流，\texttt{<=} 覆盖关系运算，\texttt{*} 和 \texttt{+} 覆盖算术运算，赋值语句覆盖变量更新。

\path{src/factorial.c} 与它的算法相同，只增加 \texttt{\#include "sysy\_runtime.h"}，并把主函数写成标准 C 形式 \texttt{int main(void)}，使 GCC 和 Clang 能检查输入输出函数声明。

\subsection{本步结论}

阶乘程序覆盖整型变量、赋值、循环条件、加法、乘法和运行库函数调用，可逐项观察其降低过程。它没有独立的 \texttt{if/else}、用户自定义函数及更多运算，因此步骤十四另用特性示例补充，而不把阶乘单独写成覆盖全部特征。

\section{步骤二：准备输入输出运行时}

SysY 使用的 \texttt{getint} 和 \texttt{putint} 不是 C 标准库函数。只编译主程序时可以生成目标文件，但链接器找不到它们的实现。因此在 \path{src/sysy_runtime.c} 中准备最小运行时：

\begin{lstlisting}[language=C]
int getint(void) {
    int value = 0;
    if (scanf("%d", &value) != 1) return 0;
    return value;
}

void putint(int value) {
    printf("%d\n", value);
}
\end{lstlisting}

\texttt{getint} 从标准输入读取一个整数，读取失败时返回 0；\texttt{putint} 输出整数和换行。C 程序和手写 LLVM IR 都连接该文件，因此两者的输入输出条件一致。本步只准备便于观察链接过程的自写最小实现，不等于已连接提供的 SysY 库。提供的库中 \texttt{putint} 没有换行，二者输出字节并不完全相同，具体在步骤十三验证。

\section{步骤三：执行预处理并观察结果}

\subsection{命令和参数}

\begin{lstlisting}[language=bash]
gcc -E -I src src/factorial.c -o build/factorial.i
\end{lstlisting}

\texttt{-E} 表示只做预处理；\texttt{-I src} 把 \path{src} 加入头文件搜索目录；最后的 \texttt{-o} 指定输出文件。检查命令为：

\begin{lstlisting}[language=bash]
wc -l build/factorial.i
sed -n '1,40p' build/factorial.i
\end{lstlisting}

实际输出共 31 行。\texttt{\#include "sysy\_runtime.h"} 已经消失，原位置出现展开后的声明：

\begin{lstlisting}[language=C]
int getint(void);
void putint(int value);
\end{lstlisting}

文件中还有 \texttt{\# 1 "src/factorial.c"} 等行标记，记录文本原来来自哪个文件和哪一行，便于后续错误定位。

\subsection{本步结论}

预处理器没有生成机器码，它主要完成头文件展开、宏替换和条件编译，输出仍是 C 文本。直接证据是 \texttt{\#include} 已消失，而声明和主函数仍可阅读。本步对应作业的“了解预处理器并获得阶段输出”。

\section{步骤四：生成未优化 LLVM IR}

\subsection{命令和参数}

\begin{lstlisting}[language=bash]
clang -S -emit-llvm -O0 -Xclang -disable-O0-optnone \
  -I src src/factorial.c -o build/factorial_clang_O0.ll
\end{lstlisting}

\texttt{-S} 输出文本形式；\texttt{-emit-llvm} 指定输出 LLVM IR；\texttt{-O0} 关闭常规优化，便于逐句对照；\texttt{-Xclang -disable-O0-optnone} 不给函数添加阻止后续优化的属性。

\subsection{变量和调用的对应}

实际 IR 中有 4 个 \texttt{alloca}、6 个 \texttt{load} 和 6 个 \texttt{store}。未优化时局部变量通常在栈上分配，再通过加载和存储访问：

\begin{lstlisting}
%5 = call i32 @getint()
store i32 %5, ptr %3, align 4   ; n = getint()
store i32 2, ptr %2, align 4    ; i = 2
store i32 1, ptr %4, align 4    ; f = 1
\end{lstlisting}

\texttt{call i32 @getint()} 调用返回 32 位整数的函数，结果为 \texttt{\%5}；\texttt{store} 再把结果写入代表 \texttt{n} 的栈地址。

\subsection{循环的对应}

\begin{lstlisting}
%9 = icmp sle i32 %7, %8
br i1 %9, label %10, label %16
%13 = mul nsw i32 %11, %12
store i32 %13, ptr %4, align 4
%15 = add nsw i32 %14, 1
store i32 %15, ptr %2, align 4
br label %6
\end{lstlisting}

\texttt{icmp sle} 对两个有符号 32 位整数做“小于等于”比较；\texttt{br i1} 进入循环体或退出块；\texttt{mul} 对应 \texttt{f*i}；\texttt{add} 对应 \texttt{i+1}；最后的无条件分支回到条件块。

\subsection{本步结论}

LLVM IR 已把结构化 \texttt{while} 变成控制流图，但仍不绑定某种 CPU 指令。变量在未优化 IR 中主要表现为栈地址，表达式结果表现为只赋值一次的 SSA 临时值。本步对应作业的“观察编译器中间输出并分析与源程序的关系”。

\section{步骤五：生成宿主机汇编}

\begin{lstlisting}[language=bash]
gcc -S -O0 -I src src/factorial.c -o build/factorial_gcc_O0.s
\end{lstlisting}

这里 \texttt{-S} 表示到汇编文本为止。未指定交叉目标，所以产物是 x86-64 汇编；它用于观察流程，最终手写目标仍选 RISC-V。关键部分为：

\begin{lstlisting}
call    getint@PLT
movl    %eax, -4(%rbp)      # n
movl    $2, -12(%rbp)       # i = 2
movl    $1, -8(%rbp)        # f = 1
.L3:
movl    -8(%rbp), %eax
imull   -12(%rbp), %eax     # f * i
movl    %eax, -8(%rbp)
addl    $1, -12(%rbp)       # i + 1
.L2:
movl    -12(%rbp), %eax
cmpl    -4(%rbp), %eax
jle     .L3
\end{lstlisting}

\texttt{eax} 保存整数返回值；\texttt{imull} 完成乘法；\texttt{cmpl} 设置条件码，\texttt{jle} 根据条件码回到循环体。同一个循环在 LLVM IR 中用 \texttt{icmp+br}，在 x86-64 中用 \texttt{cmp+jle}。这说明后端负责指令选择、寄存器使用、栈帧和跳转标签。

\section{步骤六：生成目标文件并检查符号}

\begin{lstlisting}[language=bash]
gcc -c -O0 -I src src/factorial.c -o build/factorial.o
file build/factorial.o
nm -u build/factorial.o
\end{lstlisting}

\texttt{-c} 表示完成预处理、编译和汇编，但不链接。实际结果为：

\begin{lstlisting}
build/factorial.o: ELF 64-bit LSB relocatable, x86-64
                 U getint
                 U putint
\end{lstlisting}

\texttt{relocatable} 表示它是可重定位目标文件；字母 \texttt{U} 表示两个函数在当前文件中被使用但没有定义。目标文件保存机器码、符号表和重定位信息，等待链接器解析跨文件引用。本步对应“了解汇编器生成的目标文件”。

\section{步骤七：链接并运行 C 版本}

\begin{lstlisting}[language=bash]
gcc -O0 src/factorial.c src/sysy_runtime.c -o build/factorial_c
printf '5\n' | build/factorial_c
\end{lstlisting}

GCC 分别编译主程序和运行时，再调用链接器合并两部分。主程序中的 \texttt{getint/putint} 由运行时提供，运行时中的 \texttt{scanf/printf} 由系统 C 库提供。管道左侧模拟输入 5，实际输出为 120，与 $5!=120$ 相符。

使用 \texttt{file build/factorial\_c} 可确认它是 x86-64 ELF 可执行文件，不再是可重定位目标文件。到这一步，预处理、编译、汇编和链接四个阶段已完整走通。本步对应“了解链接器并生成最终程序”。

\section{步骤八：手写并运行 LLVM IR}

\path{ir/factorial_manual.ll} 分为 \texttt{entry}、\texttt{cond}、\texttt{body}、\texttt{exit} 四个基本块：

\begin{lstlisting}
declare i32 @getint()
declare void @putint(i32)
define i32 @main() {
entry:
  %n = call i32 @getint()
  br label %cond
cond:
  %i = phi i32 [ 2, %entry ], [ %next, %body ]
  %f = phi i32 [ 1, %entry ], [ %prod, %body ]
  %cmp = icmp sle i32 %i, %n
  br i1 %cmp, label %body, label %exit
body:
  %prod = mul nsw i32 %f, %i
  %next = add nsw i32 %i, 1
  br label %cond
exit:
  call void @putint(i32 %f)
  ret i32 0
}
\end{lstlisting}

SSA 要求每个名字只定义一次，而 \texttt{i} 和 \texttt{f} 每轮都要更新。\texttt{phi} 根据控制流来源选值：第一次从 \texttt{entry} 到达时，\texttt{i=2,f=1}；以后从 \texttt{body} 返回时，取 \texttt{\%next} 和 \texttt{\%prod}。它表达的正是“上一轮结果进入下一轮”。

编译、链接和运行命令为：

\begin{lstlisting}[language=bash]
clang ir/factorial_manual.ll src/sysy_runtime.c -o build/factorial_ir
printf '5\n' | build/factorial_ir
\end{lstlisting}

Clang 把手写 IR 生成目标代码，C 运行时提供输入输出，再链接为 x86-64 程序。输出为 120，与 C 版本相同。本步直接对应“编写等价 LLVM IR 并验证”。

\section{步骤九：自动生成 RISC-V 汇编作为参考}

\begin{lstlisting}[language=bash]
clang --target=riscv64-unknown-elf -S -O0 -I src \
  src/factorial.c -o build/factorial_clang_riscv64_O0.s
\end{lstlisting}

\texttt{--target=riscv64-unknown-elf} 把目标从默认 x86-64 改为 64 位 RISC-V ELF。自动生成代码中的关键部分为：

\begin{lstlisting}
addi  sp, sp, -32
sd    ra, 24(sp)
sd    s0, 16(sp)
call  getint
sw    a0, -28(s0)       # n
li    a0, 2
sw    a0, -24(s0)       # i
li    a0, 1
sw    a0, -32(s0)       # f
lw    a1, -24(s0)
lw    a0, -28(s0)
blt   a0, a1, .LBB0_3   # n < i 时退出
\end{lstlisting}

RISC-V 调用约定使用 \texttt{a0} 保存第一个参数和整数返回值；\texttt{sp} 是栈指针，\texttt{ra} 是返回地址，\texttt{s0} 在未优化代码中作为帧指针。源条件 \texttt{i<=n} 被改写为“若 \texttt{n<i} 则退出”，因此使用 \texttt{blt}。本步验证 LLVM 的多目标后端，并给手写版本提供调用约定参考。

\section{步骤十：手写 RISC-V 阶乘主程序}

手写程序把 \texttt{n} 放在 \texttt{s0}，把 \texttt{f} 放在 \texttt{s1}，把 \texttt{i} 放在 \texttt{t0}。\texttt{getint} 返回值从 \texttt{a0} 复制到 \texttt{s0}；调用 \texttt{putint} 前把 \texttt{s1} 复制到 \texttt{a0}。

\begin{lstlisting}
main:
    addi sp, sp, -32
    sd ra, 24(sp)        # 保存返回地址
    sd s0, 16(sp)        # 保存被调用者保存寄存器
    sd s1, 8(sp)
    call getint
    mv s0, a0            # n
    li t0, 2             # i = 2
    li s1, 1             # f = 1
.Lcond:
    bgt t0, s0, .Lexit   # i > n 时退出
    mul s1, s1, t0       # f = f * i
    addi t0, t0, 1       # i = i + 1
    j .Lcond
.Lexit:
    mv a0, s1
    call putint
    li a0, 0
    ld s1, 8(sp)
    ld s0, 16(sp)
    ld ra, 24(sp)
    addi sp, sp, 32
    ret
\end{lstlisting}

\texttt{s0/s1} 属于被调用者保存寄存器，因此进入时压栈、返回前恢复；\texttt{ra} 也必须保存，因为主函数还会调用其他函数。栈帧保持 16 字节对齐。\texttt{bgt} 和 \texttt{j} 是伪指令；\texttt{mul} 来自 RISC-V M 扩展。手写版本比 \texttt{-O0} 自动代码的栈访问少，但必须人工遵守调用约定。本步对应“RISC-V 汇编编程”。

\section{步骤十一：连接 RISC-V 运行时并运行}

\subsection{运行时做了什么}

\path{asm/sysy_runtime_riscv64.s} 实现最小 SysY 兼容运行时。\texttt{\_start} 调用 \texttt{main}，随后用系统调用号 93 退出；\texttt{getint} 用系统调用号 63 从标准输入读取字符，再逐位转换十进制数；\texttt{putint} 反向生成数字和换行，再用系统调用号 64 写到标准输出。

这里的“兼容”仅指本阶乘示例所需的受限接口：输入必须从十进制数字开始，不能用它验证前导空白、正负号或多次连续读取的完整语义；输出仅处理本实验范围内的非负数。此文件也没有实现数组、浮点和计时接口，因此不称为完整 SysY 运行库。

\subsection{汇编和链接}

\begin{lstlisting}[language=bash]
riscv64-unknown-elf-as -march=rv64gc -mabi=lp64 -mno-relax \
  asm/factorial_riscv64.s -o build/factorial_riscv64.o
riscv64-unknown-elf-as -march=rv64gc -mabi=lp64 -mno-relax \
  asm/sysy_runtime_riscv64.s -o build/sysy_runtime_riscv64.o
riscv64-unknown-elf-ld --no-relax \
  build/factorial_riscv64.o build/sysy_runtime_riscv64.o \
  -o build/factorial_riscv64
\end{lstlisting}

\texttt{-march=rv64gc} 指定 64 位通用 RISC-V 指令集；\texttt{-mabi=lp64} 指定 ABI；\texttt{-mno-relax/--no-relax} 禁用链接松弛，避免最小裸入口尚未初始化 \texttt{gp} 时产生依赖 \texttt{gp} 的地址访问。

检查目标文件：

\begin{lstlisting}[language=bash]
file build/factorial_riscv64
riscv64-unknown-elf-readelf -h build/factorial_riscv64
\end{lstlisting}

实际关键信息为：

\begin{lstlisting}
ELF 64-bit LSB executable, UCB RISC-V, statically linked
Type: EXEC (Executable file)
Machine: RISC-V
Entry point address: 0x10120
\end{lstlisting}

运行命令为：

\begin{lstlisting}[language=bash]
printf '5\n' | qemu-riscv64 build/factorial_riscv64
\end{lstlisting}

实际输出 120。QEMU 在 x86-64 主机上解释执行 RISC-V 用户态指令。本步完成“汇编源文件 $\rightarrow$ 目标文件 $\rightarrow$ 连接运行时 $\rightarrow$ RISC-V ELF $\rightarrow$ 运行验证”。

\section{步骤十二：进阶观察}

\subsection{调试选项}

\begin{lstlisting}[language=bash]
clang -S -emit-llvm -g -O0 -Xclang -disable-O0-optnone \
  -I src src/factorial.c -o build/factorial_clang_g_O0.ll
\end{lstlisting}

未加 \texttt{-g} 的 IR 为 58 行，加入后为 101 行。4 个 \texttt{alloca}、6 个 \texttt{load} 和 6 个 \texttt{store} 未改变，但新增 \texttt{!dbg}、\texttt{llvm.dbg.declare}、\texttt{DICompileUnit} 等元数据，把 IR 指令、源码行号和变量关联起来。说明 \texttt{-g} 主要增加调试信息，不改变计算语义。

\subsection{优化选项}

\begin{lstlisting}[language=bash]
clang -S -emit-llvm -O2 -I src src/factorial.c \
  -o build/factorial_clang_O2.ll
\end{lstlisting}

\texttt{-O0} 通过 \texttt{alloca/load/store} 访问变量；\texttt{-O2} 中这些指令消失，并出现 9 行 \texttt{phi} 相关代码。Clang 还把部分循环向量化，出现：

\begin{lstlisting}
%11 = phi <4 x i32> [ <i32 1, i32 1, i32 1, i32 1>, %6 ], [ %15, %9 ]
%15 = mul <4 x i32> %11, %13
%22 = tail call i32 @llvm.vector.reduce.mul.v4i32(<4 x i32> %21)
\end{lstlisting}

这说明优化器能删除不必要内存访问，并把多次标量乘法组织成向量运算。形式虽然改变，但可观察输出应保持不变。

\subsection{目标架构和输入变化}

默认 GCC 生成 x86-64 的 \texttt{movl/imull/cmpl/jle}；指定 RISC-V 后生成 \texttt{lw/sw/mulw/blt/j}。算法相同，但寄存器、调用约定和指令由目标架构决定。

另外测试 0、1、5、7。0 和 1 验证循环一次都不执行时 \texttt{f=1} 是否正确，5 和 7 验证多轮循环。测试命令为：

\begin{lstlisting}[language=bash]
for n in 0 1 5 7; do
  printf '%s\n' "$n" | build/factorial_c
  printf '%s\n' "$n" | build/factorial_ir
  printf '%s\n' "$n" | qemu-riscv64 build/factorial_riscv64
done
\end{lstlisting}

本步对应 PPT 鼓励的调试、优化、目标选项和修改测试条件等进阶探索。

\section{步骤十三：提供的 SysY 运行库验证与分析}
\label{sec:provided-runtime}

\subsection{本步要解决什么问题}

前面的实验把 \texttt{getint/putint} 接口替换成自写实现，能够帮助理解符号解析，但不能作为“已直接连接提供的 SysY 库”的证据。本步使用用户提供的 \path{lib.tar.gz}，分别验证包内 x86 静态库、检查包内 RISC-V 静态库，并在必要时用包内源码重建适配当前 Linux 环境的库。以下结果来自实际运行，原二进制与重建二进制始终使用不同文件名。

\subsection{第一步：确认输入材料和文件用途}

先查看压缩包目录，再只提取本实验需要的四个文件：

\begin{lstlisting}[language=bash]
tar -tzf /home/abc/lib.tar.gz
tar -xzf /home/abc/lib.tar.gz --no-same-owner \
  lib/sylib.c lib/sylib.h lib/libsysy_x86.a lib/libsysy_riscv.a
sha256sum -c lib/SHA256SUMS
\end{lstlisting}

\path{sylib.c} 提供实现，\path{sylib.h} 提供接口和计时相关数据，两个 \path{.a} 分别是不同目标架构的静态库。压缩包还含 AArch64 库与共享库，本实验不使用，故没有引入项目。\path{lib/README.md} 记录压缩包来源与 SHA-256；四项校验均为 \texttt{OK}，说明导入后未改动这些文件。已有仓库副本时只需校验，不必反复解压。

\subsection{第二步：阅读接口，确定应该比较什么}

包内实现的关键语句为：

\begin{lstlisting}[language=C]
int getint() {
    int t;
    scanf("%d", &t);
    return t;
}
void putint(int a) { printf("%d", a); }
\end{lstlisting}

这里有三个需要解释的区别。第一，\texttt{putint(120)} 的标准输出只有字符 \texttt{120}，没有换行；原自写实现输出 \texttt{120} 加换行。第二，\texttt{scanf("\%d",...)} 支持有效整数前的空白和正号，而原 RISC-V 最小解析器不支持这两种情况。因此本步增加带前导空白和正号的输入测试，但不要求旧解析器通过这些超出其约定的测试。第三，提供的 \texttt{getint} 未检查 \texttt{scanf} 的返回值，因此本实验不通过空输入或非法字符去主张其行为有定义。

库还带有 \texttt{before\_main} 构造函数和 \texttt{after\_main} 析构函数。后者即使没有计时调用，也会向标准错误输出 \texttt{TOTAL: 0H-0M-0S-0us}。这不是阶乘的标准输出，不应混入答案比较，也不能把未进行计时的零值当作性能测量结果。

主程序继续使用 \path{src/sysy_runtime.h} 的两个函数声明，但不再链接 \path{src/sysy_runtime.c}。声明不提供函数实现，函数定义实际来自指定的库。提供的完整 \path{sylib.h} 中还定义了计时全局变量，只让 \path{sylib.c} 包含它，可以避免多个翻译单元重复定义这些对象；没有修改原头文件来掩盖问题。

\subsection{第三步：直接连接包内 x86 静态库}

先在宿主机检查接口和静态库链接是否成立：

\begin{lstlisting}[language=bash]
make sysy-host
# The target runs:
gcc -O0 src/factorial.c lib/libsysy_x86.a \
  -o build/sysy/factorial_c_host
clang ir/factorial_manual.ll lib/libsysy_x86.a \
  -o build/sysy/factorial_ir_host
make test-sysy-host
\end{lstlisting}

主程序放在库之前，使链接器处理其未定义的 \texttt{getint/putint} 后，从静态库抽取提供定义的目标文件。两个命令均未链接自写输入输出实现。实测 C 与手写 LLVM IR 共 30 次检查通过：每个程序测试 0 至 12，以及前导空白和正号两种输入格式。此结果证明包内 \path{libsysy_x86.a} 的直接链接路线成功，但不能以宿主机成功代替 RISC-V 验证。

\subsection{第四步：检查原 RISC-V 库的架构和依赖}

\begin{lstlisting}[language=bash]
make sysy-inspect
riscv64-unknown-elf-readelf -h -A lib/libsysy_riscv.a
riscv64-unknown-elf-nm -u lib/libsysy_riscv.a
\end{lstlisting}

关键结果是：

\begin{lstlisting}
File: lib/libsysy_riscv.a(sylib.o)
Class: ELF64
Type: REL (Relocatable file)
Machine: RISC-V
Flags: 0x5, RVC, double-float ABI
Tag_RISCV_stack_align: 16-bytes
U _impure_ptr
U fprintf
U gettimeofday
U printf
U putchar
U scanf
U vfprintf
\end{lstlisting}

\texttt{REL} 表示可重定位目标，静态库不是可以直接交给 QEMU 的可执行程序。\texttt{U} 表示尚需外部提供定义；因此库中虽有 SysY 函数，却仍依赖底层 C 标准库。\texttt{double-float ABI} 对应本实验采用的 \texttt{lp64d}，不能与前面 \texttt{lp64} 的目标文件不加检查地混用。ABI 除寄存器传参外还约束栈对齐，本实验继续保持 16 字节对齐\cite{riscv-abi}。即使程序不计算浮点数，目标文件的 ABI 标记也需要一致。

\subsection{第五步：保留直接链接失败的真实原因}

首先使用当前已有的裸机工具链：

\begin{lstlisting}[language=bash]
riscv64-unknown-elf-gcc -march=rv64gc -mabi=lp64d \
  asm/factorial_riscv64.s lib/libsysy_riscv.a \
  -o build/sysy/factorial_original_archive
\end{lstlisting}

该命令退出状态为 1，报告缺少 \texttt{crt0.o}、\texttt{-lc} 和 \texttt{-lgloss}。这分别涉及启动目标、C 标准库及底层支撑库，不是阶乘循环的语法错误。此前自写 \texttt{\_start} 并直接发起 Linux 系统调用的程序不需要这些文件，所以其成功不能排除本问题。

补充 RISC-V Linux 工具链后，再直接连接原库：

\begin{lstlisting}[language=bash]
riscv64-linux-gnu-gcc -march=rv64gc -mabi=lp64d -static \
  asm/factorial_riscv64.s lib/libsysy_riscv.a \
  -o build/sysy/factorial_original_linux
\end{lstlisting}

结果仍失败，错误为 \texttt{undefined reference to `\_impure\_ptr'}。Newlib 文档将此符号解释为指向其重入状态结构的全局指针\cite{newlib}，结合符号检查，可以推断原库依赖 Newlib 环境，不能与当前 glibc 环境直接混用。压缩包没有提供原构建命令，所以这里只陈述依赖证据和当前环境下的失败，不断言原库在其配套环境中损坏。

尽管主程序只调用 \texttt{getint/putint}，这两个函数与其他运行库函数位于同一个 \texttt{sylib.o} 中。本次普通静态链接抽取该成员时，其他函数的外部引用也需要解析，因此仍会遇到标准流相关的依赖。不能随意伪造一个同名全局变量来修补链接，因为名称相同不代表数据布局和运行语义相同。

\subsection{第六步：从提供的源码重建 RISC-V Linux 运行库}

采用与 QEMU 用户态环境匹配的 \texttt{riscv64-linux-gnu-gcc}，重新编译未修改的 \path{lib/sylib.c}。本次交叉 GCC 为 13.3.0，目标 glibc 为 2.39。先生成目标文件，再归档：

\begin{lstlisting}[language=bash]
riscv64-linux-gnu-gcc -march=rv64gc -mabi=lp64d -O0 \
  -c lib/sylib.c -o build/sysy/sylib_riscv_linux.o
riscv64-linux-gnu-ar rcs build/sysy/libsysy_riscv_linux.a \
  build/sysy/sylib_riscv_linux.o
\end{lstlisting}

\texttt{-c} 只生成目标文件；\texttt{ar rcs} 建立含符号索引的静态库。新库名明确包含 \texttt{linux}，没有覆盖原来的 \path{lib/libsysy_riscv.a}。这一步保留 SysY 实现源码，但改用当前目标的头文件和 C 库环境生成机器码。

用新库分别连接 C、手写 LLVM IR 和手写 RISC-V 汇编：

\begin{lstlisting}[language=bash]
riscv64-linux-gnu-gcc -march=rv64gc -mabi=lp64d -static -O0 \
  src/factorial.c build/sysy/libsysy_riscv_linux.a \
  -o build/sysy/factorial_c_riscv
clang --target=riscv64-linux-gnu -march=rv64gc -mabi=lp64d \
  -c ir/factorial_manual.ll -o build/sysy/factorial_ir_riscv.o
riscv64-linux-gnu-gcc -march=rv64gc -mabi=lp64d -static \
  build/sysy/factorial_ir_riscv.o build/sysy/libsysy_riscv_linux.a \
  -o build/sysy/factorial_ir_riscv
riscv64-linux-gnu-gcc -march=rv64gc -mabi=lp64d -static \
  asm/factorial_riscv64.s build/sysy/libsysy_riscv_linux.a \
  -o build/sysy/factorial_asm_riscv
\end{lstlisting}

\texttt{-static} 将所需库链接进可执行文件，运行时不需要另行提供目标动态加载器；但编译和链接时仍然需要 RISC-V 的头文件、启动文件和静态库。由 GCC 驱动链接可以加入匹配的启动及退出代码，使运行库构造、析构和标准 I/O 收尾正常工作。三个命令均不加入 \path{asm/sysy_runtime_riscv64.s}，避免重复定义或仍然调用自写实现。

本机没有系统级安装权限，实际将 Ubuntu 交叉工具链软件包解包到用户目录，没有修改系统安装。上述命令展示标准安装环境的写法；本机复现时先配置以下两个环境变量，并将 sysroot 传给 Makefile：

\begin{lstlisting}[language=bash]
TOOLCHAIN="$HOME/.local/share/complier-riscv-toolchain/root"
export PATH="$TOOLCHAIN/usr/bin:$PATH"
export LD_LIBRARY_PATH="$TOOLCHAIN/usr/lib/x86_64-linux-gnu${LD_LIBRARY_PATH:+:$LD_LIBRARY_PATH}"
make test-sysy RISCV_SYSROOT="$TOOLCHAIN"
\end{lstlisting}

使用用户目录工具链手动执行上面的 GCC 命令时，也需加入 \texttt{--sysroot="\$TOOLCHAIN"}。协作者若已正常安装 \texttt{gcc-riscv64-linux-gnu} 和 \texttt{libc6-dev-riscv64-cross}，则无需这段本机环境设置，直接执行 \texttt{make test-sysy} 即可。Clang 的目标三元组覆盖警告来自手写 IR 没有固定目标；本次显式指定目标并检查最终 ELF，不将该警告误判为运行失败。

\subsection{第七步：检查最终 ELF 并实际运行}

\begin{lstlisting}[language=bash]
file build/sysy/factorial_asm_riscv
riscv64-unknown-elf-readelf -h build/sysy/factorial_asm_riscv
riscv64-unknown-elf-nm -g build/sysy/factorial_asm_riscv \
  | grep -E ' (getint|putint|main|before_main|after_main)$'
printf '5\n' | qemu-riscv64 build/sysy/factorial_asm_riscv \
  > build/sysy/example.stdout 2> build/sysy/example.stderr
od -An -t x1 build/sysy/example.stdout
cat build/sysy/example.stderr
\end{lstlisting}

三个目标程序均被识别为静态链接的 RISC-V ELF，带 RVC 和双精度浮点 ABI；\texttt{readelf} 显示 \texttt{Type: EXEC}。符号表中上述五个函数均有 \texttt{T} 类型的代码定义。输入 5 后退出状态为 0，标准输出字节为 \texttt{31 32 30}，即 \texttt{120} 且无换行；标准错误为 \texttt{TOTAL: 0H-0M-0S-0us}。这与源码中的输出协议一致。

\subsection{第八步：自动测试、证据位置与结论边界}

\begin{lstlisting}[language=bash]
make test
make test-sysy
# For the local extracted toolchain, add RISCV_SYSROOT as above.
\end{lstlisting}

\path{tests/check_factorial.py} 用 Python 标准库独立计算阶乘期望值，逐次检查进程正常退出和标准输出的精确字节，并设置 5 秒超时。它不使用会吞掉尾部换行的 shell 命令替换，也不把相同的错误结果当作通过。标准错误单独保存，便于检查运行库日志，而不是丢弃。任一用例失败都会使脚本以非零状态结束，后续用例通过不会覆盖前面的失败。

实测汇总如下：

\begin{center}
\begin{tabular}{p{7.0cm}rr}
\toprule
路线 & 用例数 & 通过数 \\
\midrule
自写运行时：宿主 C、宿主 IR、RISC-V 汇编 & $13\times3$ & 39 \\
包内 x86 静态库：宿主 C、宿主 IR & $15\times2$ & 30 \\
源码重建 Linux 库：RISC-V C、IR、汇编 & $15\times3$ & 45 \\
\bottomrule
\end{tabular}
\end{center}

13 个数值用例为 0 至 12，额外两个输入分别是前导空格加制表符的 5、带正号的 5。\texttt{make test-sysy} 合并后为 \texttt{75/75 PASS}，加上基础对照共 114 次检查。每次输入、标准输出和标准错误，以及汇总 \path{results.tsv}，保存在 \path{build/test-results/baseline/} 和 \path{build/test-results/sysy/}。这些是可再生成的实验记录，不将大量生成文件提交仓库。

因此，本步完成的是：包内 x86 库直接验证成功；包内 SysY 源码重建后，RISC-V 的三种程序验证成功；原包内 RISC-V 静态库直接链接在当前环境失败，已解释并保留该限制。本步只验证了运行库的 \texttt{getint/putint} 及正常启动退出，不涉及库中数组、浮点和计时 API 的全面测试。若课程明确要求指定的 RISC-V 预编译二进制，则仍需其配套 Newlib 工具链、启动文件及执行环境。

\section{步骤十四：覆盖更多 SysY 特征并观察编译器内部阶段}
\label{sec:features}

\subsection{第一步：设计可逐项核对的程序}

仅靠阶乘无法观察独立的 \texttt{if/else}、用户函数、数组、除法和取余，因此增加 \path{src/features.sy}。相同算法的 C 文件为 \path{src/features.c}，只增加输入输出函数声明并把 \texttt{main} 写为标准 C 形式。程序定义全局整型常量 \texttt{upper=20}，从输入读取三个整数存入 \texttt{data[3]}，对范围 $0<x<20$ 内的值调用 \texttt{transform}，范围外的值加 1，最后输出和。

\begin{lstlisting}[language=C]
const int upper = 20;
int transform(int x) {
    int y;
    y = (x + 3) * 2 - x / 2;
    if (x % 2 == 0) y = y + 7;
    else y = y - 5;
    return y;
}
// main 中的核心循环：
while (i < 3) {
    data[i] = getint();
    if (data[i] > 0 && data[i] < upper)
        sum = sum + transform(data[i]);
    else
        sum = sum + 1;
    i = i + 1;
}
\end{lstlisting}

完整程序以源文件为准。\texttt{+,-,*,/ ,\%} 各有实际用途；\texttt{==} 检查奇偶，\texttt{>} 和 \texttt{<} 检查区间，\texttt{\&\&} 展示短路条件，\texttt{if/else} 分别决定加 7、减 5 或输入范围外加 1，\texttt{while} 管理三次读取，\texttt{transform} 展示形参、返回值和用户函数调用。数组下标和全局常量属于在阶乘基础上的额外观察点。程序不声称覆盖总体要求中所有一元、逻辑和控制语句、浮点或多维数组；该总体要求中的完整语言实现属于后续编译器模块。

\subsection{第二步：从预处理到词法、语法与控制流}

\begin{lstlisting}[language=bash]
make analysis
# The target runs GCC -E, Clang dump-tokens/ast-dump/DumpCFG,
# Clang -emit-llvm -O0/-O2, GCC -S/-c and RV64 -S.
\end{lstlisting}

\texttt{gcc -E} 生成 \path{build/features.i}，能看到展开的 \texttt{getint/putint} 声明和仍可阅读的函数体；它不负责判定循环分支的语义。Clang 的 \texttt{-dump-tokens} 生成 \path{build/features.tokens.txt}，本机结果 170 行，其中 \texttt{percent '\%'}、\texttt{ampamp '\&\&'}、\texttt{l\_square '['}、\texttt{if}、\texttt{while} 分别对应取余、逻辑与、数组下标、分支和循环。每个 token 还带源码位置，例如逻辑与位于 \path{src/features.c} 第 24 行；这说明词法阶段已经区分标识符、关键字和运算符。

\texttt{-ast-dump} 生成 \path{build/features.ast.txt}，本机结果 146 行。树中 \texttt{FunctionDecl transform} 明确函数及整型参数，\texttt{VarDecl data 'int[3]'} 明确数组类型，\texttt{WhileStmt}、\texttt{IfStmt} 和 \texttt{CallExpr} 体现嵌套语法关系；多个 \texttt{ArraySubscriptExpr} 分别对应数组写入、条件读取与传给函数。AST 还能显示表达式分层，而 token 列表不能单独表示括号优先级或分支层次。

\texttt{DumpCFG} 生成 \path{build/features.cfg.txt}。在 \texttt{main} 的控制流图中，\texttt{while} 条件块有进入循环和退出两条后继边；\texttt{\&\&} 被拆成两次判断，第一个条件不成立就跳到范围外分支，不执行第二次比较。\texttt{transform} 的奇偶分支最后汇合到同一个返回块。这些实际块和边使结构化控制语句如何降为跳转可以被直接检查。

\subsection{第三步：比较自动 LLVM IR 与手写 SSA}

自动生成的 \path{build/features_clang_O0.ll} 保留多处 \texttt{alloca/load/store}，便于回找局部变量。\texttt{getelementptr inbounds [3 x i32]} 计算 \texttt{data[i]} 的元素地址，\texttt{sdiv} 和 \texttt{srem} 分别实现有符号除法、取余；对 \texttt{transform} 的调用由 \texttt{call i32 @transform} 表示。生成的 \texttt{-O2} IR 结构明显不同，部分条件和调用已被优化，不应用源代码行数直接衡量运行性能，因本实验没有做计时基准。

手写 \path{ir/features_manual.ll} 中的 \texttt{transform} 分成 \texttt{even\_path} 和 \texttt{odd\_path}，两条路径的计算结果由 \texttt{phi} 合并。\texttt{main} 的 \texttt{cond} 块用两个 \texttt{phi} 接收上轮的 \texttt{i} 和 \texttt{sum}。\texttt{data} 用 \texttt{alloca [3 x i32]} 分配；\texttt{getelementptr} 计算元素地址后 \texttt{store/load}，明确表现数组访问。判断 \texttt{data[i]>0} 后才进入 \texttt{upper\_check}，因此第二个比较只在第一个条件为真时执行，符合短路语义。

\begin{lstlisting}
%i = phi i32 [ 0, %entry ], [ %next, %join ]
%sum = phi i32 [ 0, %entry ], [ %new_sum, %join ]
%slot = getelementptr inbounds [3 x i32], ptr %data, i32 0, i32 %i
%changed = call i32 @transform(i32 %value)
\end{lstlisting}

两种 IR 不需要指令逐字相同，只需在本程序定义的输入范围内具有一致的可观察行为。手写 IR 使用有符号 32 位运算；测试值较小，避免有符号溢出造成的语义分歧。

\subsection{第四步：解释手写 RISC-V 的数组和函数调用}

手写 \path{asm/features_riscv64.s} 在 \texttt{main} 建立 64 字节栈帧：前 12 字节为三个整型元素，栈中另外保存 \texttt{ra/s0/s1}。\texttt{s0} 表示循环下标，\texttt{s1} 表示累计和，二者在函数调用后仍需保留，所以按照调用约定压栈并在返回前恢复。\texttt{slli t0,s0,2} 把元素下标乘以 4，再与 \texttt{sp} 相加计算 \texttt{data[i]} 的地址；\texttt{sw/lw} 实现 32 位写入和读取。

\texttt{blez} 先处理 $x\leq0$，\texttt{bge} 再处理 $x\geq20$，两次分支对应短路条件。满足范围时将数组元素放入 \texttt{a0} 调用 \texttt{transform}。函数内部使用 \texttt{divw/remw} 完成有符号 32 位除法和取余，并用 \texttt{addiw/subw} 保持 32 位结果。调用 \texttt{putint} 前再把 \texttt{s1} 放入 \texttt{a0}。工具链将它与步骤十三重建的 Linux 运行库链接，\texttt{file} 确认所得程序为静态 RISC-V ELF；QEMU 用户态能执行。

\subsection{第五步：构建、核对输出与有效范围}

\begin{lstlisting}[language=bash]
make analysis
make test-features
# User-local toolchain: make test-features RISCV_SYSROOT="$TOOLCHAIN"
\end{lstlisting}

\texttt{test-features} 分别构建宿主 C、宿主手写 IR、RISC-V C、RISC-V 手写 IR、RISC-V 手写汇编，共五个程序。它们只连接提供的 x86 库或根据提供的 \texttt{sylib.c} 重建的 RISC-V Linux 库。测试脚本独立计算期望值、核对进程退出状态及标准输出的完整字节，并保存每组输入、输出和标准错误到 \path{build/test-results/features/}。

\begin{center}
\begin{tabular}{p{4cm}rp{7cm}}
\toprule
输入三个整数 & 期望和 & 主要覆盖 \\
\midrule
\texttt{2 3 4} & 41 & 偶数、奇数，两次分支与函数调用 \\
\texttt{0 20 -3} & 3 & 零、上界、负数都走范围外分支 \\
\texttt{19 1 18} & 73 & 上界以内，数组三次写入与读取 \\
\texttt{7 7 7} & 36 & 重复输入、多次循环与累计 \\
\bottomrule
\end{tabular}
\end{center}

另外还测试 \texttt{0 0 0}、\texttt{1 2 19}、\texttt{20 21 22}、\texttt{-1 0 1}。实测八组输入乘以五个程序为 \texttt{40/40 PASS}；不是只检查不同程序互相一致，而是各自对照期望值。提供的 \texttt{putint} 无换行，因此比较字节时不应期望尾部换行。该程序只读三个有效整数；非法输入、整数溢出、浮点和多维数组仍不在此结论范围。

\section{实验结果与分析}

\begin{center}
\begin{tabular}{ccccc}
\toprule
输入 $n$ & 期望值 & C 版本 & 手写 LLVM IR & 手写 RISC-V \\
\midrule
0 & 1 & 1 & 1 & 1 \\
1 & 1 & 1 & 1 & 1 \\
5 & 120 & 120 & 120 & 120 \\
7 & 5040 & 5040 & 5040 & 5040 \\
\bottomrule
\end{tabular}
\end{center}

上表为原先四组测试，分别覆盖循环不执行和多轮执行。现有自动测试进一步覆盖 $0\leq n\leq12$ 的全部整数，三种自写运行时版本共 39 次检查全部通过。期望值由测试脚本独立计算，不仅仅互相比对三个程序，避免相同错误被误判为正确。

提供的 SysY 库路线还有 75 次检查，详见步骤十三；新多特征示例有 40 次检查，详见步骤十四。三条路线合计 154 次检查全部通过；其中不同示例、运行库与目标组合各自有不同的用例集合，不把数量当作互相独立的数学证明。

C 和 LLVM IR 使用 32 位整数；当前手写 RV64 版本使用 64 位 \texttt{mul/addi}，不能声称其溢出行为与前两者相同。$12!=479001600$ 可由有符号 32 位整数表示，$13!=6227020800$ 已超出范围，因此只在 $0\leq n\leq12$ 的约定有效域内比较。负数不是本阶乘实验的有效输入；无效输入、EOF 和整数溢出未纳入正确性结论。有限测试是证据，不是普遍证明。

\section{遇到的问题与处理}

第一，只生成 \path{build/factorial.o} 时，\texttt{getint/putint} 仍是未定义符号，文件不能直接运行。增加运行时并执行链接后解决，这也说明了汇编器和链接器职责的区别。

第二，默认 GCC 输出 x86-64，不能证明 RISC-V 代码正确。因此使用 Clang 目标三元组生成 RISC-V 参考汇编，再用 RISC-V GNU 工具链处理手写代码，通过 \texttt{readelf} 和 QEMU 验证目标架构及结果。

第三，RISC-V 链接器可能把地址加载放松为相对 \texttt{gp} 的访问，但本实验的最小 \texttt{\_start} 没有 C 启动文件初始化 \texttt{gp}。因此汇编时使用 \texttt{-mno-relax}，链接时使用 \texttt{--no-relax}，运行时也声明 \texttt{.option norelax}。

\section{每一步与作业要求的对应关系}

\begin{enumerate}
  \item \textbf{步骤一、二与十四}对应“设计 SysY 示例并覆盖主要语言特征”。阶乘便于跟踪编译流程；补充示例加入全局常量、整型数组、加减乘除取余、条件分支、逻辑与、循环和用户函数。未覆盖总体要求中的所有语法，更未实现一个完整编译器。
  \item \textbf{步骤三}对应“了解预处理器并获得预处理输出”。证据是 \texttt{gcc -E} 和 \path{build/factorial.i}。
  \item \textbf{步骤四、五}对应“了解编译器并分析源程序到 IR、汇编的关系”。证据是未优化 LLVM IR、x86-64 汇编及逐句对照。
  \item \textbf{步骤六}对应“了解汇编器和目标文件”。证据是 \path{build/factorial.o} 的 ELF 类型、未定义符号和可重定位状态。
  \item \textbf{步骤七}对应“了解链接器并生成最终程序”。证据是主程序与 C 运行时链接，以及输入 5 输出 120。
  \item \textbf{步骤八}对应“编写等价 LLVM IR 并验证”。证据是 \path{ir/factorial_manual.ll}、SSA/\texttt{phi} 分析和 \path{build/factorial_ir}。
  \item \textbf{步骤九、十、十一}对应“ARM 或 RISC-V 任选其一，连接运行时并生成目标程序验证”的自写最小运行时路线。本实验选择 RISC-V，证据是参考汇编、手写汇编、ELF 和 QEMU 输出；提供的 SysY 库另外在步骤十三验证。
  \item \textbf{步骤十二}对应进阶要求。完成调试信息比较、\texttt{-O0/-O2} 比较、x86-64/RISC-V 目标比较和多输入测试。
  \item \textbf{步骤十三}补充“连接 SysY 运行库并运行”的直接证据，以及目标文件 ABI、标准库依赖和输出协议的进阶分析。具体完成边界以该步骤的实测结果为准。
  \item \textbf{步骤十四}对应 PPT 中“更细致的编译器各阶段功能”和“涵盖语言特征的等价 IR/汇编”，证据是 token、AST、CFG、自动 IR、手写 IR、手写 RISC-V 汇编以及五种可执行程序的 40 次检查。数组为额外探索，不将它当作已经完成后续编译器模块的进阶语法实现。
  \item \textbf{摘要、关键词、任务说明、实验过程、结果、结论、参考文献及分工}对应“按科技论文结构撰写 PDF 报告”。本文件是待个人补充、审核和成功编译的报告源稿；列出 PPT 模板链接不等于已经使用其版式，也不等于已经形成可提交的 PDF。
\end{enumerate}

\section{结论}

本实验完整观察了预处理、编译、汇编和链接过程，并用真实中间文件解释各阶段职责。源程序的 \texttt{while} 在 LLVM IR 中变成基本块、比较和分支，在 RISC-V 中变成寄存器比较、条件分支和回跳；局部变量在未优化 IR 中通过内存读写表示，在手写 SSA IR 中通过 \texttt{phi} 合并控制流来源；函数调用则需要在链接阶段由运行时提供定义。

在此基础上，本实验手写两个示例的 LLVM IR 和 RISC-V 主程序，在约定输入域内获得一致结果。第二个示例还让 token、AST、CFG、数组元素地址及短路分支之间的对应关系能够直接观察。对 \texttt{-g}、\texttt{-O0}、\texttt{-O2} 和目标架构选项的比较展示了中间产物的形式变化；这里的优化部分主要是静态观察，不把没有执行的性能基准或所有优化级别测试算作完成。提供的 SysY 库验证还表明，目标指令集正确并不足以保证链接成功，必须同时考虑调用约定、C 标准库、启动流程和输入输出协议。

\begin{thebibliography}{99}
\bibitem{course} 课程材料：《预备工作——了解你的编译器》PPT 与《上机大作业总体要求》文档。
\bibitem{llvm-langref} LLVM Project. LLVM Language Reference Manual. \url{https://llvm.org/docs/LangRef.html}
\bibitem{clang} LLVM Project. Clang Command Line Argument Reference. \url{https://clang.llvm.org/docs/ClangCommandLineReference.html}
\bibitem{gcc} GNU Project. GCC Online Documentation. \url{https://gcc.gnu.org/onlinedocs/}
\bibitem{riscv} RISC-V International. The RISC-V Instruction Set Manual. \url{https://riscv.org/technical/specifications/}
\bibitem{riscv-abi} RISC-V International. RISC-V ABIs Specification. \url{https://riscv-non-isa.github.io/riscv-elf-psabi-doc/}
\bibitem{newlib} Newlib. The Newlib C Library, Reentrancy. \url{https://sourceware.org/newlib/libc.html}
\bibitem{overleaf-template} 课程 PPT 推荐报告模板：\url{https://www.overleaf.com/read/hnsjbjyknbdj}
\bibitem{latex-intro} LaTeX 入门资料：\url{https://liam.page/2014/09/08/latex-introduction/}
\bibitem{latex-symbols} LaTeX 命令与符号资料：\url{https://blog.csdn.net/garfielder007/article/details/51646604}
\bibitem{latex-equation} LaTeX 数学公式资料：\url{https://yundongxiaoyang.top/wiki/latex-equation/}
\end{thebibliography}

\end{document}
